\section*{Isotropic elements and singular elements:}

Let \(\Phi\) be a form \(\varepsilon\)-hermitian on a left A-module E. An element \(x \in E\) is said to be isotropic if \(\Phi(x, x) = 0\) ; a submodule M of E is said to be isotropic if \(\mathbf{M} \cap \mathbf{M}^{\circ} \neq \{0\}\) (in other words, if the restriction of \(\Phi\) to M is degenerate); M is said to be totally isotropic if \(\mathbf{M} \subset \mathbf{M}^{\circ}\) (in other words, if the restriction of \(\Phi\) to M is zero). For every submodule M of E, \(\mathbf{M} \cap \mathbf{M}^{\circ}\) is totally isotropic. If E is a finite-dimensional vector space and if \(\Phi\) is nondegenerate, the following three conditions are equivalent for a subspace M of E : \(1^{\circ}\) M is nonisotropic ; \(2^{\circ}\) M \(^{\circ}\) is nonisotropic ; \(3^{\circ}\) E is the direct sum of M and M \(^{\circ}\) . Suppose the ring A is commutative, and let Q be a quadratic form on an A-module E, \(\Phi\) the bilinear form associated with Q. An element \(x \in E\) is said to be singular if \(\mathrm{Q}(x) = 0\) ; a submodule M of E is said to be singular (resp. totally singular) if \(\mathbf{M} \cap \mathbf{M}^{\circ}\) contains a singular element \(\neq 0\) (resp. if the restriction of Q to M is zero). Every singular element (resp. singular submodule, totally singular submodule) is isotropic (resp. isotropic, totally isotropic); the converse is true when A is a field of characteristic \(\neq 2\) .
